\documentclass[11pt, a4paper]{article}\usepackage[utf8]{inputenc}\usepackage{amsmath}\usepackage[margin=1in]{geometry}\begin{document}\title{Problem Set 10: Konsep GA}\maketitle
\section*{Soal 1 (C2)} Apa dampak dari laju mutasi (mutation rate) yang terlalu tinggi terhadap pencarian solusi pada GA?
\section*{Soal 2 (C3)} Jika kita memiliki kromosom biner: P1=[1,0,1,0,1] dan P2=[0,1,1,1,0]. Tunjukkan hasil Crossover satu-titik (One-point crossover) jika titik potongnya berada setelah gen kedua!
\end{document}